% !TEX TS-program = pdflatex
\documentclass[11pt,a4paper]{article}
\usepackage[utf8]{inputenc}
\usepackage[T1]{fontenc}
\usepackage[default]{lato}
\usepackage[margin=2.5cm]{geometry}
\usepackage{xcolor}
\usepackage[hidelinks]{hyperref}
\usepackage{parskip} % for spacing between paragraphs instead of indent

\definecolor{Navy}{HTML}{0A2540}

\begin{document}

\begin{flushleft}
\textbf{\Large Achraf IZELMADENE}\\
Marrakech, Maroc (Mobilité vers Montréal)\\
+212 699 098 736\\
\href{mailto:achraf.izelmadene@gmail.com}{achraf.izelmadene@gmail.com}\\
\href{https://www.linkedin.com/in/achraf-izelmadene/}{linkedin.com/in/achraf-izelmadene}
\end{flushleft}

\vspace{0.5cm}

\begin{flushright}
Le 15 août 2026
\end{flushright}

\vspace{0.3cm}

\textbf{Lantic Inc.}\\
4026 Notre-Dame Est\\
Montréal, QC\\
Canada

\vspace{0.5cm}

\textbf{Objet : Candidature au poste de Stagiaire – Génie Électrique / Génie de la Production Automatisée (Projet LEAP)}

\vspace{0.5cm}

Madame, Monsieur,

C'est avec un grand intérêt que je vous soumets ma candidature pour le stage au sein de l'équipe du projet d'expansion LEAP. Actuellement étudiant ingénieur en quatrième année (plus de 60 crédits accumulés) en Génie Électrique et Systèmes Intelligents à l'EMSI, je suis passionné par l'automatisation industrielle et l'instrumentation, ce qui fait de ce stage chez Lantic une opportunité idéale pour moi.

Lors de mon stage en unité de production chez Menara Prefa, j'ai participé activement à la migration d'un automate vers un PLC Siemens S7 (TIA Portal). Cette expérience m'a permis de réaliser des vérifications I/O, d'effectuer des tests fonctionnels sur des équipements de production et de collaborer à la résolution de déficiences techniques sur le terrain. J'y ai développé une forte débrouillardise et une capacité d'apprentissage rapide dans des environnements industriels exigeants.

Par la suite, mon passage en bureau d'études chez Core Engineering Solutions m'a familiarisé avec la lecture, l'interprétation et la mise à jour de plans électriques, de schémas unifilaires et de P\&ID. En participant au dimensionnement HTA/BT, j'ai appris l'importance d'assurer une parfaite concordance entre les installations sur le terrain et la documentation technique. 

De plus, mes projets académiques, tels que la modélisation de boucles de contrôle sous Simulink ou la conception de réseaux de capteurs (instrumentation) pour une agriculture intelligente, témoignent de mon fort intérêt pour les variateurs de vitesse, les essais d'équipements et les systèmes de contrôle (HMI/SCADA).

Participer concrètement à la mise en service d'une nouvelle usine et assurer le suivi des installations électriques au sein du projet LEAP représente pour moi le défi technique parfait. Je suis convaincu que ma curiosité, ma rigueur et ma double compétence en génie électrique et production automatisée me permettront de supporter efficacement votre équipe multidisciplinaire.

Je vous remercie par avance du temps que vous consacrerez à l'étude de ma candidature. Je me tiens à votre entière disposition pour une entrevue afin de discuter de ma motivation.

Veuillez agréer, Madame, Monsieur, l'expression de mes salutations distinguées.

\vspace{1cm}
\textbf{Achraf IZELMADENE}

\end{document}
